\chapter{Chapter06 — Low Degree F₂Polynomial Learning}
%%% generated by the blueprint pipeline from FABL.Chapter06.LearningAndTesting.LowDegreeF₂PolynomialLearning
\section{Overview}

Declarations of the Lean module \texttt{FABL.Chapter06.LearningAndTesting.LowDegreeF₂PolynomialLearning}.

\section{Declarations}

\begin{definition}[Linear equation over the binary field]
\lean{FABL.F₂LinearEquation}
\leanok
\uses{FABL.𝔽₂}
For a type $\nu$, an $\mathbb{F}_2$-linear equation consists of a coefficient function
$a : \nu \to \mathbb{F}_2$ and a constant $c \in \mathbb{F}_2$.
\end{definition}

\begin{definition}[Evaluation of a binary linear equation]
\lean{FABL.F₂LinearEquation.eval}
\leanok
\uses{FABL.F₂LinearEquation, FABL.𝔽₂}
Let $V$ be a finite set, let $a_v	ext{ and }x_v$ be elements of $br/2br$ for each
$v	ext{ in }V$, and let $E$ be the equation with coefficients $a_v$. Its evaluation at
the assignment $x$ is $sum_{v\in V} a_vx_v$ in $br/2br$.
\end{definition}

\begin{definition}[Satisfaction of a binary linear equation]
\lean{FABL.F₂LinearEquation.IsSatisfied}
\leanok
\uses{FABL.F₂LinearEquation, FABL.F₂LinearEquation.eval, FABL.𝔽₂}
Let $V$ be a finite set, let $E$ be a linear equation over $br/2br$ with constant term
$c$, and let $x:V\to\bbr/2\bbr$ be an assignment. The assignment satisfies $E$ exactly
when the sum of each coefficient times its assigned value equals $c$ in $br/2\bbr$.
\end{definition}

\begin{definition}[Addition of binary linear equations]
\lean{FABL.F₂LinearEquation.add}
\leanok
\uses{FABL.F₂LinearEquation}
Let $V$ be a finite set, and let $E_1$ and $E_2$ be linear equations over $br/2br$ on
$V$. Their sum is the equation whose coefficient at each variable is the sum of the
corresponding coefficients of $E_1$ and $E_2$, and whose constant term is the sum of
their constant terms, all in $br/2br$.
\end{definition}

\begin{theorem}[Evaluation preserves equation addition]
\lean{FABL.F₂LinearEquation.eval_add}
\leanok
\uses{FABL.F₂LinearEquation, FABL.F₂LinearEquation.add, FABL.F₂LinearEquation.eval, FABL.𝔽₂}
\difficulty{1}
Let $\nu$ be a finite set, let $L_1,L_2$ be linear equations over $\mathbb{F}_2$
on $\nu$, and let $x\colon\nu\to\mathbb{F}_2$ be an assignment. The evaluation of the
sum of the two equations at $x$ equals the sum of their evaluations at $x$.
\end{theorem}

\begin{theorem}[Satisfaction under addition of a satisfied equation]
\lean{FABL.F₂LinearEquation.isSatisfied_add_iff_of_isSatisfied_right}
\leanok
\uses{FABL.F₂LinearEquation, FABL.F₂LinearEquation.IsSatisfied, FABL.F₂LinearEquation.add, FABL.F₂LinearEquation.eval, FABL.F₂LinearEquation.eval_add, FABL.𝔽₂}
\difficulty{4}
Let $\nu$ be a finite set, let $L_1,L_2$ be linear equations over $\mathbb{F}_2$ on
$\nu$, and let $x\colon\nu\to\mathbb{F}_2$ be an assignment satisfying $L_2$. Then $x$
satisfies the sum of $L_1$ and $L_2$ if and only if it satisfies $L_1$.
\end{theorem}

\begin{definition}[Clearing a pivot variable]
\lean{FABL.F₂LinearEquation.clearPivot}
\leanok
\uses{FABL.F₂LinearEquation, FABL.F₂LinearEquation.add}
Let $V$ be a finite set, choose a variable $p\in V$, and let $P$ and $R$ be linear
equations over $rr/2rr$ on $V$. The result of clearing the pivot from $R$ is $R$
itself if its coefficient at $p$ is zero, and otherwise is the sum $R+P$.
\end{definition}

\begin{theorem}[Pivot coefficient after clearing]
\lean{FABL.F₂LinearEquation.clearPivot_coefficient_eq_zero}
\leanok
\uses{FABL.F₂LinearEquation, FABL.F₂LinearEquation.clearPivot}
\difficulty{3}
Let a pivot variable and two linear equations over $\mathbb{F}_2$ on a finite
set of variables be given. If the pivot equation has coefficient $1$ at the pivot
variable, then after clearing that pivot from the row equation, the resulting equation
has coefficient $0$ at the pivot variable.
\end{theorem}

\begin{theorem}[Satisfaction after clearing a pivot]
\lean{FABL.F₂LinearEquation.isSatisfied_clearPivot_iff}
\leanok
\uses{FABL.F₂LinearEquation, FABL.F₂LinearEquation.IsSatisfied, FABL.F₂LinearEquation.clearPivot, FABL.F₂LinearEquation.isSatisfied_add_iff_of_isSatisfied_right, FABL.𝔽₂}
\difficulty{3}
Let $\nu$ be a finite set, let $p$ be a pivot variable, and let $P$ and $R$ be linear
equations over $\mathbb{F}_2$ on $\nu$. For any assignment $x\colon\nu\to\mathbb{F}_2$
satisfying $P$, the cleared row equation is satisfied by $x$ if and only if $R$ is
satisfied by $x$.
\end{theorem}

\begin{definition}[Assignments satisfying a list of equations]
\lean{FABL.F₂SatisfiesRows}
\leanok
\uses{FABL.F₂LinearEquation, FABL.F₂LinearEquation.IsSatisfied, FABL.𝔽₂}
Let $V$ be a finite set, and let each row in a finite list of equations assign a
coefficient in the field $br/2br$ to every element of $V$ and have a constant in that
field. An assignment $x:V\to\bbr/2\bbr$ satisfies the list if it satisfies every
equation that occurs in it.
\end{definition}

\begin{definition}[Extracting a pivot equation]
\lean{FABL.extractF₂Pivot}
\leanok
\uses{FABL.F₂LinearEquation}
Let $V$ be a finite set, choose a variable $p\in V$, and let $mathcal R$ be a list of
linear equations over $rr/2rr$ on $V$. Pivot extraction returns no result if no row
has coefficient $1$ at $p$; otherwise it returns the first such row together with the
list of all other rows in their original order.
\end{definition}

\begin{theorem}[Coefficient of an extracted pivot]
\lean{FABL.extractF₂Pivot_coefficient_eq_one}
\leanok
\uses{FABL.F₂LinearEquation, FABL.extractF₂Pivot}
\difficulty{4}
Let $\nu$ be a finite set, choose a variable $v\in\nu$, and consider a list of linear
equations over $\mathbb F_2$, each with a coefficient for every variable. If extracting
a pivot equation at $v$ from the list returns an equation $E$ and a remaining list, then
the coefficient of $v$ in $E$ is $1$.
\end{theorem}

\begin{theorem}[Satisfaction after pivot extraction]
\lean{FABL.extractF₂Pivot_satisfiesRows_iff}
\leanok
\uses{FABL.F₂LinearEquation, FABL.F₂LinearEquation.IsSatisfied, FABL.F₂SatisfiesRows, FABL.extractF₂Pivot, FABL.𝔽₂}
\difficulty{5}
Let $\nu$ be a finite set, choose a variable $v\in\nu$, and let a list of linear
equations over $\mathbb F_2$ yield a pivot equation $E$ and a remaining list under pivot
extraction at $v$. For any assignment of a value in $\mathbb F_2$ to each variable, all
equations in the original list are satisfied if and only if $E$ and every equation in
the remaining list are satisfied.
\end{theorem}

\begin{theorem}[When pivot extraction fails]
\lean{FABL.extractF₂Pivot_eq_none_iff}
\leanok
\uses{FABL.F₂LinearEquation, FABL.extractF₂Pivot}
\difficulty{5}
Let $\nu$ be a finite set, choose a variable $v\in\nu$, and let a list of linear
equations over $\mathbb F_2$ assign a coefficient to each variable. Pivot extraction at
$v$ returns no equation if and only if every equation in the list has coefficient $0$ at
$v$.
\end{theorem}

\begin{definition}[Clearing a Pivot from a List of Equations]
\lean{FABL.clearF₂PivotRows}
\leanok
\uses{FABL.F₂LinearEquation, FABL.F₂LinearEquation.clearPivot}
Let $V$ be a finite set, let $p\in V$, and let $P$ be a linear equation in variables
indexed by $V$ over the field with two elements. For any finite list of linear equations
in these variables, the result of clearing the pivot $p$ from the list is the list
obtained by applying the pivot-clearing operation with pivot $p$ and equation $P$ to
each row, in the original order.
\end{definition}

\begin{theorem}[Pivot Clearing Preserves Satisfaction]
\lean{FABL.satisfiesRows_clearF₂PivotRows_iff}
\leanok
\uses{FABL.F₂LinearEquation, FABL.F₂LinearEquation.IsSatisfied, FABL.F₂LinearEquation.clearPivot, FABL.F₂LinearEquation.isSatisfied_clearPivot_iff, FABL.F₂SatisfiesRows, FABL.clearF₂PivotRows, FABL.𝔽₂}
\difficulty{3}
Let $\nu$ be a finite set, let $p$ be a linear equation over $mathbb{F}_2$, and let $R$
be a finite list of such equations. Fix an assignment that satisfies $p$. Clearing the
coefficient of a chosen variable in every equation of $R$ by adding $p$ whenever that
coefficient is nonzero preserves satisfaction of the list: the assignment satisfies the
resulting list if and only if it satisfies $R$.
\end{theorem}

\begin{definition}[Binary elimination result]
\lean{FABL.F₂EliminationResult}
\leanok
\uses{FABL.F₂LinearEquation}
For a finite variable set $V$, an elimination result consists of a list of
pivot-variable and equation pairs, a list of residual equations, and a natural number
recording the amount of work.
\end{definition}

\begin{definition}[Recursive binary row elimination]
\lean{FABL.eliminateF₂Rows}
\leanok
\uses{FABL.F₂EliminationResult, FABL.F₂LinearEquation, FABL.clearF₂PivotRows, FABL.extractF₂Pivot}
For a finite type of variables, the procedure takes an ordered list of variables and a
list of linear equations over the two-element field. It recursively processes the
variables, using the pivot-extraction and row-clearing operations at each step, and
returns the resulting pivot list, residual rows, and a recursively accumulated
natural-number work count.
\end{definition}

\begin{theorem}[Length after pivot extraction]
\lean{FABL.extractF₂Pivot_remaining_length}
\leanok
\uses{FABL.F₂LinearEquation, FABL.extractF₂Pivot}
\difficulty{4}
Let $\nu$ be a finite set, choose a variable $v\in\nu$, and let a list of linear
equations over $\mathbb F_2$ yield a pivot equation and a remaining list under pivot
extraction at $v$. The length of the remaining list plus one equals the length of the
original list.
\end{theorem}

\begin{theorem}[Length preservation under pivot clearing]
\lean{FABL.length_clearF₂PivotRows}
\leanok
\uses{FABL.F₂LinearEquation, FABL.clearF₂PivotRows}
\difficulty{1}
Let $\nu$ be a finite set of variables, let a pivot variable and a pivot equation over
$\mathbb{F}_2$ be given, and let $R$ be a finite list of linear equations over
$\mathbb{F}_2$. Clearing the pivot from each equation in $R$ preserves the length of the
list.
\end{theorem}

\begin{theorem}[Bound on the number of pivots]
\lean{FABL.eliminateF₂Rows_pivots_length_le}
\leanok
\uses{FABL.F₂LinearEquation, FABL.clearF₂PivotRows, FABL.eliminateF₂Rows, FABL.extractF₂Pivot}
\difficulty{4}
Let $\nu$ be a finite set, and let $C$ be a list of elements of $\nu$ and $R$ a list of
linear equations over the field $\mathbb{F}_2$. The elimination procedure applied to $C$
and $R$ produces a pivot list whose length is at most the length of $C$.
\end{theorem}

\begin{theorem}[Bound on elimination work]
\lean{FABL.eliminateF₂Rows_work_le}
\leanok
\uses{FABL.F₂LinearEquation, FABL.clearF₂PivotRows, FABL.eliminateF₂Rows, FABL.extractF₂Pivot, FABL.extractF₂Pivot_remaining_length}
\difficulty{5}
Let $
u$ be a finite coordinate set, let $C$ be a list of coordinates, and let $R$ be a list
of linear equations over $mathbb{F}_2$. The work recorded by eliminating along $C$ is
at most $|C|\,|R|\,(|\nu|+2)$.
\end{theorem}

\begin{theorem}[Satisfaction After Linear Elimination]
\lean{FABL.satisfiesRows_eliminateF₂Rows_iff}
\leanok
\uses{FABL.F₂LinearEquation, FABL.F₂LinearEquation.IsSatisfied, FABL.F₂SatisfiesRows, FABL.clearF₂PivotRows, FABL.eliminateF₂Rows, FABL.extractF₂Pivot, FABL.extractF₂Pivot_satisfiesRows_iff, FABL.satisfiesRows_clearF₂PivotRows_iff, FABL.𝔽₂}
\difficulty{5}
Let $\nu$ be a finite variable set, let $C$ be a list of variables, and let $R$ be a
list of linear equations over $mathbb{F}_2$. For any assignment, it satisfies every
equation in $R$ if and only if it satisfies every pivot equation recorded by eliminating
along $C$ and every equation in the resulting residual list.
\end{theorem}

\begin{theorem}[Predicates on extracted equations]
\lean{FABL.extractF₂Pivot_forall_iff}
\leanok
\uses{FABL.F₂LinearEquation, FABL.extractF₂Pivot}
\difficulty{4}
Let $\nu$ be a finite set, choose a variable $v\in\nu$, and let a list of linear
equations over $\mathbb F_2$ yield a pivot equation $E$ and a remaining list under pivot
extraction at $v$. For any predicate $P$ on such equations, every equation in the
original list satisfies $P$ if and only if $E$ satisfies $P$ and every equation in the
remaining list satisfies $P$.
\end{theorem}

\begin{theorem}[Preservation of a zero column]
\lean{FABL.clearF₂PivotRows_column_eq_zero}
\leanok
\uses{FABL.F₂LinearEquation, FABL.F₂LinearEquation.add, FABL.F₂LinearEquation.clearPivot, FABL.clearF₂PivotRows}
\difficulty{3}
Let $\mathbb F_2$ be the field with two elements, let $\nu$ be a finite set of
variables, and choose a fixed variable, a pivot variable, a pivot equation, and a finite
list of equations. If the pivot equation has coefficient zero at the fixed variable and
every equation in the list also has coefficient zero there, then after clearing the
pivot variable from every equation in the list, every resulting equation still has
coefficient zero at the fixed variable.
\end{theorem}

\begin{theorem}[Preservation of a zero column]
\lean{FABL.eliminateF₂Rows_preserves_zero_column}
\leanok
\uses{FABL.F₂LinearEquation, FABL.clearF₂PivotRows, FABL.clearF₂PivotRows_column_eq_zero, FABL.eliminateF₂Rows, FABL.extractF₂Pivot, FABL.extractF₂Pivot_forall_iff}
\difficulty{5}
Let $
u$ be a finite coordinate set, let $z\\in\\nu$, and let $C$ be a list of coordinates and
$R$ a list of linear equations over $\\mathbb{F}_2$. If every equation in $R$ has
coefficient zero at $z$, then every pivot equation and every residual equation produced
by eliminating along $C$ also has coefficient zero at $z$.
\end{theorem}

\begin{definition}[Triangular pivot list]
\lean{FABL.IsF₂Triangular}
\leanok
\uses{FABL.F₂LinearEquation}
A list of variable–equation pairs is triangular if, for each pair in the list, its
equation has coefficient $1$ at its associated variable, every later equation has
coefficient $0$ at that variable, and the tail of the list is itself triangular.
\end{definition}

\begin{theorem}[Elimination produces triangular pivots]
\lean{FABL.isF₂Triangular_eliminateF₂Rows}
\leanok
\uses{FABL.F₂LinearEquation, FABL.F₂LinearEquation.clearPivot_coefficient_eq_zero, FABL.IsF₂Triangular, FABL.clearF₂PivotRows, FABL.eliminateF₂Rows, FABL.eliminateF₂Rows_preserves_zero_column, FABL.extractF₂Pivot, FABL.extractF₂Pivot_coefficient_eq_one}
\difficulty{5}
Let $V$ be a finite set of variables, let $c_1,\ldots,c_k$ be an ordered list of
variables in $V$, and let a finite list of linear equations over $\mathbb F_2$ in the
variables $V$ be given. Applying the specified row elimination procedure to these
variables and equations produces pivots that are triangular: each pivot equation has
coefficient $1$ on its pivot variable, every later pivot equation has coefficient $0$ on
that earlier pivot variable, and these properties hold at every position in the pivot
list.
\end{theorem}

\begin{definition}[Set a pivot value]
\lean{FABL.setF₂PivotValue}
\leanok
\uses{FABL.F₂LinearEquation, FABL.𝔽₂}
Given a finite variable type, a pivot variable, a linear equation, and an assignment of
field values to all variables, form a new assignment by changing only the pivot
variable's value. Set that value to the equation's constant plus the sum, over all other
variables, of the coefficient at that variable times its value in the original
assignment.
\end{definition}

\begin{theorem}[Pivot Update Preserves Other Equations]
\lean{FABL.isSatisfied_setF₂PivotValue_iff_of_coefficient_eq_zero}
\leanok
\uses{FABL.F₂LinearEquation, FABL.F₂LinearEquation.IsSatisfied, FABL.F₂LinearEquation.eval, FABL.setF₂PivotValue, FABL.𝔽₂}
\difficulty{4}
Let $\\nu$ be a finite set of variables, let $E$ and $R$ be equations over $\\mathbb
F_2$ on $\\nu$, and let $x:\\nu\\to\\mathbb F_2$ be an assignment. If the coefficient of
the pivot variable in $R$ is $0$, then $R$ is satisfied by the assignment obtained by
setting the pivot value using $E$ if and only if $R$ is satisfied by $x$.
\end{theorem}

\begin{theorem}[Satisfaction After Solving for a Pivot]
\lean{FABL.isSatisfied_setF₂PivotValue}
\leanok
\uses{FABL.F₂LinearEquation, FABL.F₂LinearEquation.IsSatisfied, FABL.F₂LinearEquation.eval, FABL.setF₂PivotValue, FABL.𝔽₂}
\difficulty{4}
Let $\\nu$ be a finite set of variables, let an equation over $\\mathbb F_2$ have
coefficients $a_v$ and constant $c$, and let $x:\\nu\\to\\mathbb F_2$ be an assignment.
If the coefficient of a chosen pivot variable $p$ is $1$, replace $x(p)$ by
$c+\\sum_{v\\ne p}a_vx(v)$, leaving the other values unchanged. The resulting assignment
satisfies the equation.
\end{theorem}

\begin{definition}[Solve an ordered pivot list]
\lean{FABL.solveF₂Pivots}
\leanok
\uses{FABL.F₂LinearEquation, FABL.setF₂PivotValue, FABL.𝔽₂}
For a finite variable type, recursively turn a list of variable–equation pairs into an
assignment: the empty list gives the identically zero assignment, and a nonempty list
applies the pivot-value update for its first pair to the assignment obtained from the
remaining list.
\end{definition}

\begin{theorem}[Satisfaction of triangular equations]
\lean{FABL.solveF₂Pivots_satisfies}
\leanok
\uses{FABL.F₂LinearEquation, FABL.F₂LinearEquation.IsSatisfied, FABL.IsF₂Triangular, FABL.isSatisfied_setF₂PivotValue, FABL.isSatisfied_setF₂PivotValue_iff_of_coefficient_eq_zero, FABL.solveF₂Pivots}
\difficulty{4}
Let $\nu$ be a finite set and let a list of pairs $(v_i,E_i)$ consist of variables
$v_i\in\nu$ and linear equations over $\mathbb F_2$. Suppose each equation has
coefficient $1$ at its paired variable, and every later equation in the list has
coefficient $0$ at each earlier paired variable. The assignment obtained by setting the
paired variables recursively to satisfy their equations satisfies every equation in the
list.
\end{theorem}

\begin{theorem}[Residual coefficients vanish at processed coordinates]
\lean{FABL.eliminateF₂Rows_residual_coefficient_eq_zero_of_mem}
\leanok
\uses{FABL.F₂LinearEquation, FABL.F₂LinearEquation.clearPivot_coefficient_eq_zero, FABL.clearF₂PivotRows, FABL.eliminateF₂Rows, FABL.eliminateF₂Rows_preserves_zero_column, FABL.extractF₂Pivot, FABL.extractF₂Pivot_coefficient_eq_one, FABL.extractF₂Pivot_eq_none_iff}
\difficulty{5}
Let $
u$ be a finite coordinate set, let $C$ be a list of coordinates, and let $R$ be a list
of linear equations over $mathbb{F}_2$. For every coordinate $z$ that occurs in $C$,
every residual equation produced by eliminating along $C$ has coefficient zero at $z$.
\end{theorem}

\begin{definition}[Sorted coordinate list]
\lean{FABL.f₂LinearCoordinateList}
\leanok
For a finite encodable type, the coordinate list is its list of elements in the order
given by the sorted enumeration of its universal finite set.
\end{definition}

\begin{theorem}[Membership in the coordinate list]
\lean{FABL.mem_f₂LinearCoordinateList}
\leanok
\uses{FABL.f₂LinearCoordinateList}
\difficulty{1}
Let $\nu$ be a finite, encodable type. Every coordinate $x\nu$ belongs to its sorted
list of all elements.
\end{theorem}

\begin{theorem}[Length of the linear coordinate list]
\lean{FABL.length_f₂LinearCoordinateList}
\leanok
\uses{FABL.f₂LinearCoordinateList}
\difficulty{1}
For every finite, encodable type of coordinates, the linear coordinate list has length
equal to the number of elements of that type.
\end{theorem}

\begin{definition}[Solve a list of binary equations]
\lean{FABL.solveF₂Rows}
\leanok
\uses{FABL.F₂LinearEquation, FABL.eliminateF₂Rows, FABL.f₂LinearCoordinateList, FABL.solveF₂Pivots, FABL.𝔽₂}
For a finite encodable type with decidable equality, the procedure first eliminates the
given list of linear equations using the sorted coordinate list, then returns the
assignment obtained by solving the resulting pivot list.
\end{definition}

\begin{theorem}[Consistency of the linear system solver]
\lean{FABL.solveF₂Rows_satisfies_of_exists}
\leanok
\uses{FABL.F₂LinearEquation, FABL.F₂LinearEquation.IsSatisfied, FABL.F₂LinearEquation.eval, FABL.F₂SatisfiesRows, FABL.eliminateF₂Rows, FABL.eliminateF₂Rows_residual_coefficient_eq_zero_of_mem, FABL.f₂LinearCoordinateList, FABL.isF₂Triangular_eliminateF₂Rows, FABL.satisfiesRows_eliminateF₂Rows_iff, FABL.solveF₂Pivots_satisfies, FABL.solveF₂Rows, FABL.𝔽₂}
\difficulty{5}
Let $\nu$ be a finite set equipped with an encoding, and let $R$ be a finite list of
linear equations over $\mathbb F_2$ in variables indexed by $\nu$. If some assignment
satisfies every equation in $R$, then the assignment produced by the specified solver
also satisfies every equation in $R$.
\end{theorem}

\begin{definition}[Low degree monomials]
\lean{FABL.LowDegreeMonomial}
\statementsource{Analysis-of-Boolean-Functions-by-Ryan-ODonnell}{
  pdf-ffbe9c7b3306-p149-b001,
  pdf-ffbe9c7b3306-p149-b004,
  pdf-ffbe9c7b3306-p167-b005
}
\label{FABL.LowDegreeMonomial}
\leanok
For natural numbers $n$ and $\ell$, the low degree monomials are indexed by subsets
$S\subseteq\{0,\ldots,n-1\}$ with $|S|\leq\ell$.
\end{definition}

\begin{definition}[Low degree binary coefficients]
\lean{FABL.LowDegreeF₂Coefficients}
\leanok
\uses{FABL.LowDegreeMonomial, FABL.𝔽₂}
\statementsource{Analysis-of-Boolean-Functions-by-Ryan-ODonnell}{
  pdf-ffbe9c7b3306-p149-b001,
  pdf-ffbe9c7b3306-p149-b004,
  pdf-ffbe9c7b3306-p167-b005
}
For natural numbers $n$ and $\ell$, a family of low degree binary coefficients assigns
an element of $\mathbb F_2$ to each subset of $\{0,\ldots,n-1\}$ of size at most $\ell$.
\end{definition}

\begin{definition}[Low degree monomial count]
\lean{FABL.lowDegreeF₂MonomialCount}
\leanok
\uses{FABL.LowDegreeMonomial}
For natural numbers $n$ and $\ell$, the low-degree monomial count is the cardinality of
the type of low-degree monomials indexed by $n$ and bounded by degree $\ell$.
\end{definition}

\begin{theorem}[Counting low degree monomials]
\lean{FABL.lowDegreeF₂MonomialCount_eq_sum_choose}
\leanok
\uses{FABL.card_lowDegreeFourierFamily_eq_sum_choose, FABL.lowDegreeFourierFamily, FABL.lowDegreeF₂MonomialCount}
\difficulty{4}
For every $n,\ell\in\mathbb{N}$, the number of subsets of $[n]$ of size at most $\ell$
is $\sum_{j=0}^{\ell}\binom{n}{j}$.
\end{theorem}

\begin{theorem}[Bound on the number of low degree monomials]
\lean{FABL.lowDegreeF₂MonomialCount_le}
\leanok
\uses{FABL.card_lowDegreeFourierFamily_eq_sum_choose, FABL.card_lowDegreeFourierFamily_le, FABL.lowDegreeF₂MonomialCount, FABL.lowDegreeF₂MonomialCount_eq_sum_choose}
\difficulty{3}
For every $n,\ell\in\mathbb{N}$, the number of subsets of $[n]$ of size at most $\ell$
is at most $(\ell+1)(n+1)^\ell$.
\end{theorem}

\begin{definition}[Evaluation of a low degree binary polynomial]
\lean{FABL.lowDegreeF₂Eval}
\leanok
\uses{FABL.F₂Cube, FABL.LowDegreeF₂Coefficients, FABL.anfMonomial, FABL.𝔽₂}
\statementsource{Analysis-of-Boolean-Functions-by-Ryan-ODonnell}{
  pdf-ffbe9c7b3306-p149-b001,
  pdf-ffbe9c7b3306-p149-b004,
  pdf-ffbe9c7b3306-p167-b005
}
Let $n$ and $\ell$ be natural numbers, let $c_S\in\mathbb F_2$ be a coefficient for each
subset $S\subseteq\{0,\ldots,n-1\}$ of size at most $\ell$, and let $x\in\mathbb F_2^n$.
The evaluation is $\sum_{|S|\leq\ell} c_S m_S(x)$, where $m_S(x)$ is the binary monomial
associated with $S$.
\end{definition}

\begin{definition}[Extend low degree coefficients]
\lean{FABL.extendLowDegreeF₂Coefficients}
\leanok
\uses{FABL.ANFCoefficients, FABL.LowDegreeF₂Coefficients}
Given natural numbers $n$ and $\ell$, extend a coefficient function on subsets of
$\{0,\ldots,n-1\}$ of size at most $\ell$ to all subsets of $\{0,\ldots,n-1\}$ by
retaining the given coefficient on subsets of size at most $\ell$ and assigning $0$ to
every larger subset.
\end{definition}

\begin{theorem}[Extending low-degree coefficients]
\lean{FABL.anfEval_extendLowDegreeF₂Coefficients}
\leanok
\uses{FABL.LowDegreeF₂Coefficients, FABL.LowDegreeMonomial, FABL.anfEval, FABL.anfMonomial, FABL.extendLowDegreeF₂Coefficients, FABL.lowDegreeFourierFamily, FABL.lowDegreeF₂Eval}
\difficulty{4}
Let $\ell$ be a nonnegative integer and assign a coefficient in $\mathbb{F}_2$ to every
subset of $\{1,\ldots,n\}$ of size at most $\ell$. Extend this family by assigning
coefficient zero to every larger subset. The resulting algebraic normal form agrees, at
every $x\in\mathbb{F}_2^n$, with the sum of the assigned coefficients times their
corresponding monomials over subsets of size at most $\ell$.
\end{theorem}

\begin{theorem}[Uniqueness of low degree polynomial coefficients]
\lean{FABL.lowDegreeF₂Eval_injective}
\leanok
\uses{FABL.F₂BooleanFunction, FABL.LowDegreeF₂Coefficients, FABL.anfEval_extendLowDegreeF₂Coefficients, FABL.anfEval_injective, FABL.extendLowDegreeF₂Coefficients, FABL.lowDegreeF₂Eval}
\difficulty{4}
For every $n,
\ell\in\mathbb{N}$, the map sending a coefficient assignment on subsets of $[n]$ of size
at most $\ell$ to the Boolean function on $\mathbb{F}_2^n$ obtained by summing its
coefficients times the corresponding squarefree monomials is injective.
\end{theorem}

\begin{theorem}[Additivity of low degree evaluation]
\lean{FABL.lowDegreeF₂Eval_add}
\leanok
\uses{FABL.LowDegreeF₂Coefficients, FABL.lowDegreeF₂Eval}
\difficulty{3}
For all nonnegative integers $n$ and $\ell$, evaluation of a sum of two coefficient
families on the $n$-dimensional Boolean cube over $\mathbb{F}_2$, using monomials of
degree at most $\ell$, equals the sum of their evaluations at every input.
\end{theorem}

\begin{theorem}[Coefficients of a low-degree polynomial]
\lean{FABL.anfCoeff_lowDegreeF₂Eval}
\leanok
\uses{FABL.LowDegreeF₂Coefficients, FABL.anfCoeff, FABL.anfEval_anfCoeff, FABL.anfEval_extendLowDegreeF₂Coefficients, FABL.anfEval_injective, FABL.extendLowDegreeF₂Coefficients, FABL.lowDegreeF₂Eval}
\difficulty{2}
Let $n$ and $\ell$ be natural numbers, and assign a coefficient in $\mathbb F_2$ to each
subset of $\{0,\ldots,n-1\}$ of size at most $\ell$. Form the Boolean function on the
$n$-dimensional cube by summing each assigned coefficient times the corresponding
coordinate monomial. Its algebraic normal form coefficient at each subset is the
assigned coefficient when the subset has size at most $\ell$, and is $0$ otherwise.
\end{theorem}

\begin{theorem}[Degree of a low-degree polynomial]
\lean{FABL.functionAlgebraicDegree_lowDegreeF₂Eval_le}
\leanok
\uses{FABL.LowDegreeF₂Coefficients, FABL.algebraicDegree_le_iff, FABL.anfCoeff_lowDegreeF₂Eval, FABL.extendLowDegreeF₂Coefficients, FABL.functionAlgebraicDegree, FABL.lowDegreeF₂Eval}
\difficulty{3}
Let $n$ and $\ell$ be nonnegative integers, and assign a coefficient in $\mathbb{F}_2$
to each subset of $\{1,\ldots,n\}$ of size at most $\ell$. The function on
$\mathbb{F}_2^n$ obtained by summing these coefficients times their coordinate monomials
has algebraic degree at most $\ell$.
\end{theorem}

\begin{definition}[Coefficients of a function]
\lean{FABL.lowDegreeF₂CoefficientsOfFunction}
\leanok
\uses{FABL.F₂BooleanFunction, FABL.LowDegreeF₂Coefficients, FABL.anfCoeff}
\statementsource{Analysis-of-Boolean-Functions-by-Ryan-ODonnell}{pdf-ffbe9c7b3306-p167-b005}
Let $n$ and $\ell$ be natural numbers and let $f:\mathbb F_2^n\to\mathbb F_2$. Associate
to each subset $S\subseteq\{0,\ldots,n-1\}$ of size at most $\ell$ the coefficient of
the monomial indexed by $S$ in the algebraic normal form of $f$.
\end{definition}

\begin{theorem}[Evaluation of a function’s low degree coefficients]
\lean{FABL.lowDegreeF₂Eval_coefficientsOfFunction}
\leanok
\uses{FABL.F₂BooleanFunction, FABL.LowDegreeMonomial, FABL.algebraicDegree_le_iff, FABL.anfCoeff, FABL.anfEval, FABL.anfEval_anfCoeff, FABL.anfMonomial, FABL.functionAlgebraicDegree, FABL.lowDegreeFourierFamily, FABL.lowDegreeF₂CoefficientsOfFunction, FABL.lowDegreeF₂Eval}
\difficulty{5}
Let $n$ and $\ell$ be nonnegative integers, and let $f$ be a function from the
$n$-dimensional Boolean cube over $\mathbb{F}_2$ to $\mathbb{F}_2$. If the algebraic
degree of $f$ is at most $\ell$, then evaluating its algebraic normal form coefficients
on monomials of degree at most $\ell$ recovers $f$ at every input.
\end{theorem}

\begin{definition}[Sample equation]
\lean{FABL.lowDegreeF₂SampleEquation}
\leanok
\uses{FABL.F₂Cube, FABL.F₂LinearEquation, FABL.LowDegreeMonomial, FABL.anfMonomial, FABL.𝔽₂}
\statementsource{Analysis-of-Boolean-Functions-by-Ryan-ODonnell}{pdf-ffbe9c7b3306-p167-b005}
Let $n$ and $\ell$ be natural numbers, and let $(x,b)\in\mathbb F_2^n\times\mathbb F_2$.
The sample equation has coefficient $m_S(x)$ at each subset $S\subseteq\{0,\ldots,n-1\}$
of size at most $\ell$, and constant term $b$.
\end{definition}

\begin{theorem}[Sample equation satisfaction criterion]
\lean{FABL.lowDegreeF₂SampleEquation_isSatisfied_iff}
\leanok
\uses{FABL.F₂Cube, FABL.F₂LinearEquation.IsSatisfied, FABL.F₂LinearEquation.eval, FABL.LowDegreeF₂Coefficients, FABL.lowDegreeF₂Eval, FABL.lowDegreeF₂SampleEquation, FABL.𝔽₂}
\difficulty{3}
For any natural numbers $n,\ell$, labeled sample $(x,b)\in\mathbb F_2^n\times\mathbb
F_2$, and coefficient family on subsets of size at most $\ell$, the linear equation
formed from the sample is satisfied by that family if and only if its low-degree
polynomial evaluates at $x$ to $b$.
\end{theorem}

\begin{definition}[Rows from samples]
\lean{FABL.lowDegreeF₂SampleRows}
\leanok
\uses{FABL.F₂Cube, FABL.F₂LinearEquation, FABL.LowDegreeMonomial, FABL.lowDegreeF₂SampleEquation, FABL.𝔽₂}
\statementsource{Analysis-of-Boolean-Functions-by-Ryan-ODonnell}{pdf-ffbe9c7b3306-p167-b005}
Let $n$, $\ell$, and $m$ be natural numbers, and let $(x_i,b_i)\in\mathbb
F_2^n\times\mathbb F_2$ be given for each $i=0,\ldots,m-1$. The sample rows are the list
of $m$ linear equations whose row for $(x_i,b_i)$ has coefficient $m_S(x_i)$ at each
subset $S$ of size at most $\ell$ and constant term $b_i$.
\end{definition}

\begin{theorem}[Function coefficients satisfy sample equations]
\lean{FABL.coefficientsOfFunction_satisfies_lowDegreeF₂SampleRows}
\leanok
\uses{FABL.F₂BooleanFunction, FABL.F₂Cube, FABL.F₂SatisfiesRows, FABL.functionAlgebraicDegree, FABL.lowDegreeF₂CoefficientsOfFunction, FABL.lowDegreeF₂Eval_coefficientsOfFunction, FABL.lowDegreeF₂SampleEquation_isSatisfied_iff, FABL.lowDegreeF₂SampleRows}
\difficulty{4}
Let $f:\mathbb F_2^n\to\mathbb F_2$ be a Boolean function whose algebraic normal form
has degree at most $\ell$, and choose any finite list $x_1,\ldots,x_m\in\mathbb F_2^n$.
Assign to each monomial indexed by a set $S$ with $|S|\le\ell$ the algebraic normal form
coefficient of $f$ at $S$. This assignment satisfies every sample equation: at each
$x_i$, the sum over such $S$ of that assigned coefficient times $\prod_{q\in S}(x_i)_q$
equals $f(x_i)$ in $\mathbb F_2$.
\end{theorem}

\begin{definition}[Separation of low-degree coefficients]
\lean{FABL.SeparatesLowDegreeF₂Coefficients}
\leanok
\uses{FABL.F₂Cube, FABL.LowDegreeF₂Coefficients, FABL.lowDegreeF₂Eval}
\statementsource{Analysis-of-Boolean-Functions-by-Ryan-ODonnell}{pdf-ffbe9c7b3306-p167-b005}
Let $n$, $\ell$, and $m$ be nonnegative integers, and let $x_1,\dots,x_m$ be points of
the Boolean cube $\mathbb F_2^n$. The samples separate the low-degree coefficient arrays
if the vector of evaluations at these samples determines the array uniquely: any two
arrays indexed by subsets of $[n]$ of size at most $\ell$ that agree on all $m$
evaluations are equal.
\end{definition}

\begin{theorem}[A separation criterion for low-degree coefficients]
\lean{FABL.separatesLowDegreeF₂Coefficients_iff}
\leanok
\uses{FABL.F₂Cube, FABL.LowDegreeF₂Coefficients, FABL.SeparatesLowDegreeF₂Coefficients, FABL.lowDegreeF₂Eval, FABL.lowDegreeF₂Eval_add}
\difficulty{6}
Let $n,\ell,m$ be natural numbers, and choose $m$ inputs from the Boolean cube $\mathbb
F_2^n$. Evaluation on these inputs is injective on the coefficient arrays indexed by
subsets of $[n]$ of size at most $\ell$ if and only if every nonzero such coefficient
array has a nonzero polynomial evaluation at at least one of the chosen inputs, where
the polynomial is the sum of its coefficients times the corresponding square-free
monomials over $\mathbb F_2$.
\end{theorem}

\begin{theorem}[Uniform event probability as an indicator average]
\lean{FABL.measure_uniformPMF_event_eq_uniformProbability}
\label{FABL.measure_uniformPMF_event_eq_uniformProbability}
\leanok
\uses{FABL.integral_uniformPMF_eq_expect, FABL.uniformPMF, FABL.uniformProbability}
\difficulty{3}
Let $\Omega$ be a finite nonempty measurable space whose singletons are measurable, and
let $E\subseteq\Omega$. The uniform probability measure of $E$ equals the uniform
average over $\Omega$ of the indicator of membership in $E$.
\end{theorem}

\begin{theorem}[Uniform probability of a complement]
\lean{FABL.uniformProbability_not}
\label{FABL.uniformProbability_not}
\leanok
\uses{FABL.uniformProbability}
\difficulty{4}
Let $\Omega$ be a finite nonempty set, and let $P$ be a predicate on $\Omega$. The
uniform probability that $P$ does not hold is one minus the uniform probability that $P$
holds.
\end{theorem}

\begin{definition}[Samples on which a coefficient array vanishes]
\lean{FABL.lowDegreeF₂MissSet}
\leanok
\uses{FABL.F₂Cube, FABL.LowDegreeF₂Coefficients, FABL.lowDegreeF₂Eval}
\statementsource{Analysis-of-Boolean-Functions-by-Ryan-ODonnell}{pdf-ffbe9c7b3306-p167-b005}
Let $n$, $\ell$, and $m$ be nonnegative integers, and let $c$ be a low-degree
coefficient array indexed by subsets of $[n]$ of size at most $\ell$. The associated
miss set consists of all sequences of $m$ points of $\mathbb F_2^n$ at which the
evaluation of $c$ is zero at every sample.
\end{definition}

\begin{theorem}[Probability of missing a polynomial on the Boolean cube]
\lean{FABL.measure_lowDegreeF₂MissSet_eq_pow}
\leanok
\uses{FABL.F₂Cube, FABL.LowDegreeF₂Coefficients, FABL.lowDegreeF₂Eval, FABL.lowDegreeF₂MissSet, FABL.measure_uniformPMF_event_eq_uniformProbability, FABL.uniformPMF, FABL.uniformProbability, FABL.uniformSample_toMeasure_eq_pi}
\difficulty{3}
Let $n,\\ell,m$ be natural numbers, and let $c$ assign a coefficient in $\\mathbb F_2$
to each subset of $[n]$ of size at most $\\ell$. Form the polynomial on $\\mathbb F_2^n$
whose monomials are the products of the coordinates indexed by those subsets, with
coefficients given by $c$. For a uniformly random $m$-tuple of points of $\\mathbb
F_2^n$, the probability that this polynomial vanishes at every sampled point equals the
$m$th power of the fraction of points in $\\mathbb F_2^n$ at which it vanishes.
\end{theorem}

\begin{theorem}[Miss probability for a nonzero polynomial]
\lean{FABL.measure_lowDegreeF₂MissSet_le_exp}
\leanok
\uses{FABL.F₂Cube, FABL.LowDegreeF₂Coefficients, FABL.functionAlgebraicDegree_lowDegreeF₂Eval_le, FABL.inv_two_pow_le_uniformProbability_ne_zero_of_functionAlgebraicDegree_le, FABL.lowDegreeF₂Eval, FABL.lowDegreeF₂Eval_injective, FABL.lowDegreeF₂MissSet, FABL.measure_lowDegreeF₂MissSet_eq_pow, FABL.uniformPMF, FABL.uniformProbability, FABL.uniformProbability_not}
\difficulty{5}
Let $n,\ell,m\in\mathbb{N}$, and let $P$ be a nonzero polynomial over $\mathbb{F}_2$
expressed as an algebraic normal form using monomials indexed by subsets of $[n]$ of
size at most $\ell$. For $m$ independent uniformly sampled points of $\mathbb{F}_2^n$,
the probability that $P$ vanishes at every sampled point is at most
$\exp\!\left(-m(1/2)^\ell\right)$.
\end{theorem}

\begin{definition}[Nonzero low-degree coefficient arrays]
\lean{FABL.NonzeroLowDegreeF₂Coefficients}
\leanok
\uses{FABL.LowDegreeF₂Coefficients}
\statementsource{Analysis-of-Boolean-Functions-by-Ryan-ODonnell}{pdf-ffbe9c7b3306-p167-b005}
For nonnegative integers $n$ and $\ell$, a nonzero low-degree coefficient array is a
function from the subsets of $[n]$ of size at most $\ell$ to $\mathbb F_2$ that is not
identically zero.
\end{definition}

\begin{definition}[Failure of coefficient separation]
\lean{FABL.lowDegreeF₂SeparationFailureSet}
\leanok
\uses{FABL.F₂Cube, FABL.SeparatesLowDegreeF₂Coefficients}
\statementsource{Analysis-of-Boolean-Functions-by-Ryan-ODonnell}{pdf-ffbe9c7b3306-p167-b005}
Let $n$, $\ell$, and $m$ be nonnegative integers. The separation failure set consists of
the sequences of $m$ points of $\mathbb F_2^n$ for which the sample evaluations do not
determine every low-degree coefficient array uniquely.
\end{definition}

\begin{theorem}[Separation failure is a union of miss sets]
\lean{FABL.lowDegreeF₂SeparationFailureSet_subset_iUnion_missSet}
\leanok
\uses{FABL.NonzeroLowDegreeF₂Coefficients, FABL.lowDegreeF₂MissSet, FABL.lowDegreeF₂SeparationFailureSet, FABL.separatesLowDegreeF₂Coefficients_iff}
\difficulty{4}
For any natural numbers $n,\ell,m$, every list of $m$ inputs in $\mathbb F_2^n$ that
fails to distinguish all coefficient families on monomials of degree at most $\ell$ lies
in the miss set of some nonzero such coefficient family; that miss set consists of lists
on which the corresponding polynomial evaluates to zero at every input.
\end{theorem}

\begin{theorem}[Counting low degree binary coefficients]
\lean{FABL.card_lowDegreeF₂Coefficients}
\leanok
\uses{FABL.LowDegreeF₂Coefficients, FABL.lowDegreeF₂MonomialCount}
\difficulty{1}
For natural numbers $n$ and $\ell$, let $m$ be the number of subsets of an $n$-element
set having size at most $\ell$. The set of all functions from those subsets to the
two-element field has cardinality $2^m$.
\end{theorem}

\begin{theorem}[Measure low degree f₂separation failure set le]
\lean{FABL.measure_lowDegreeF₂SeparationFailureSet_le}
\leanok
\uses{FABL.F₂Cube, FABL.LowDegreeF₂Coefficients, FABL.NonzeroLowDegreeF₂Coefficients, FABL.card_lowDegreeF₂Coefficients, FABL.lowDegreeF₂MissSet, FABL.lowDegreeF₂MonomialCount, FABL.lowDegreeF₂SeparationFailureSet, FABL.lowDegreeF₂SeparationFailureSet_subset_iUnion_missSet, FABL.measure_lowDegreeF₂MissSet_le_exp, FABL.uniformPMF}
\statementsource{Analysis-of-Boolean-Functions-by-Ryan-ODonnell}{pdf-ffbe9c7b3306-p167-b005}
\difficulty{5}
% dataset-pipeline: stub -- prose pending build_informal_statements.py
Statement pending regeneration.
\end{theorem}

\begin{theorem}[Failure probability for polynomial separation]
\lean{FABL.measure_lowDegreeF₂SeparationFailureSet_le_of_budget}
\leanok
\uses{FABL.F₂Cube, FABL.lowDegreeF₂MonomialCount, FABL.lowDegreeF₂SeparationFailureSet, FABL.measure_lowDegreeF₂SeparationFailureSet_le, FABL.uniformPMF}
\difficulty{1}
Let $n,\ell,m\in\mathbb{N}$, and let $M$ be the number of subsets of $[n]$ of size at
most $\ell$. If $\delta\in\mathbb{R}$ satisfies
$2^M\exp\!\left(-m(1/2)^\ell\right)\le\delta$, then for $m$ independent uniform samples
from $\mathbb{F}_2^n$, the probability that evaluation on the samples fails to
distinguish some two distinct degree-at-most-$\ell$ polynomials over $\mathbb{F}_2$ is
at most $\delta$.
\end{theorem}

\begin{definition}[Low-degree learning sample count]
\lean{FABL.lowDegreeF₂LearningSampleCount}
\leanok
\uses{FABL.PositiveLearningParameter, FABL.fourierEstimatorFailureBits, FABL.lowDegreeF₂MonomialCount}
\statementsource{Analysis-of-Boolean-Functions-by-Ryan-ODonnell}{pdf-ffbe9c7b3306-p167-b005}
Let $n$ and $\ell$ be nonnegative integers, and let the rational learning parameter
$\delta$ satisfy $0<\delta\leq \tfrac12$. The prescribed sample count is $2^\ell$ times
the sum of the number of degree-at-most-$\ell$ monomials on $n$ coordinates and the
number of failure bits associated with $\delta$.
\end{definition}

\begin{theorem}[Scaled learning sample count identity]
\lean{FABL.lowDegreeF₂LearningSampleCount_mul_invTwoPow}
\leanok
\uses{FABL.PositiveLearningParameter, FABL.fourierEstimatorFailureBits, FABL.lowDegreeF₂LearningSampleCount, FABL.lowDegreeF₂MonomialCount}
\difficulty{3}
For every $n,\ell\in\mathbb{N}$ and every rational learning parameter $\delta$ with
$0<\delta\leq\frac12$, the learning sample count multiplied by $2^{-\ell}$ equals the
number of subsets of $[n]$ of size at most $\ell$ plus
$\lceil\log_2\lceil2/\delta\rceil\rceil$.
\end{theorem}

\begin{theorem}[Learning sample count budget]
\lean{FABL.lowDegreeF₂LearningSampleCount_budget_le}
\leanok
\uses{FABL.PositiveLearningParameter, FABL.fourierEstimatorFailureBits, FABL.lowDegreeF₂LearningSampleCount, FABL.lowDegreeF₂LearningSampleCount_mul_invTwoPow, FABL.lowDegreeF₂MonomialCount, FABL.two_div_pow_fourierEstimatorFailureBits_le}
\difficulty{5}
For every $n,\ell\in\mathbb{N}$ and every rational learning parameter $\delta$ with
$0<\delta\leq \frac12$, let $M$ be the number $2^\ell$ times the sum of the number of
subsets of $[n]$ of size at most $\ell$ and $\lceil\log_2\lceil2/\delta\rceil\rceil$.
Then $2^M{}^{\!0}$ is not the asserted quantity; explicitly, $2^{\#}\exp\!\left(-M
2^{-\ell}\right)\leq\delta$, where $\# $ denotes the number of subsets of $[n]$ of size
at most $\ell$.
\end{theorem}

\begin{theorem}[Measure low degree f₂separation failure set scheduled le]
\lean{FABL.measure_lowDegreeF₂SeparationFailureSet_scheduled_le}
\leanok
\uses{FABL.F₂Cube, FABL.PositiveLearningParameter, FABL.lowDegreeF₂LearningSampleCount, FABL.lowDegreeF₂LearningSampleCount_budget_le, FABL.lowDegreeF₂SeparationFailureSet, FABL.measure_lowDegreeF₂SeparationFailureSet_le_of_budget, FABL.uniformPMF}
\statementsource{Analysis-of-Boolean-Functions-by-Ryan-ODonnell}{pdf-ffbe9c7b3306-p167-b005}
\difficulty{1}
% dataset-pipeline: stub -- prose pending build_informal_statements.py
Statement pending regeneration.
\end{theorem}

\begin{theorem}[Upper bound for the learning sample count]
\lean{FABL.lowDegreeF₂LearningSampleCount_le}
\leanok
\uses{FABL.PositiveLearningParameter, FABL.fourierEstimatorFailureBits, FABL.lowDegreeF₂LearningSampleCount, FABL.lowDegreeF₂MonomialCount_le}
\difficulty{2}
For every $n,\ell\in\mathbb{N}$ and every rational learning parameter $\delta$ with
$0<\delta\leq\frac12$, the sample count $2^\ell$ times the sum of the number of subsets
of $[n]$ of size at most $\ell$ and $\lceil\log_2\lceil2/\delta\rceil\rceil$ is at most
$2^\ell\big((\ell+1)(n+1)^\ell+\lceil\log_2\lceil2/\delta\rceil\rceil\big)$.
\end{theorem}

\begin{definition}[Coordinatewise sign-to-bit encoding]
\lean{FABL.signSamplesF₂Equiv}
\leanok
\uses{FABL.F₂Cube, FABL.binaryCubeSignEquiv}
For natural numbers $n$ and $m$, the map sending each of $m$ samples of sign vectors of
length $n$ to the corresponding samples of binary vectors is the coordinatewise
bijection induced by the standard bijection between signs and elements of $bF_2$.
\end{definition}

\begin{definition}[Sign-coded separation failure set]
\lean{FABL.signSampleLowDegreeF₂SeparationFailureSet}
\leanok
\uses{FABL.lowDegreeF₂SeparationFailureSet, FABL.signSamplesF₂Equiv}
\statementsource{Analysis-of-Boolean-Functions-by-Ryan-ODonnell}{pdf-ffbe9c7b3306-p167-b005}
Let $n$, $\ell$, and $m$ be nonnegative integers. The sign-coded separation failure set
consists of those sequences of $m$ vectors in $\{−1,1\}^n$ whose corresponding
Boolean-cube samples fail to separate the low-degree coefficient arrays.
\end{definition}

\begin{theorem}[Scheduled low degree coefficient separation]
\lean{FABL.measure_signSampleLowDegreeF₂SeparationFailureSet_scheduled_le}
\leanok
\uses{FABL.F₂Cube, FABL.PositiveLearningParameter, FABL.lowDegreeF₂LearningSampleCount, FABL.lowDegreeF₂SeparationFailureSet, FABL.map_uniformPMF_equiv, FABL.measure_lowDegreeF₂SeparationFailureSet_scheduled_le, FABL.signSampleLowDegreeF₂SeparationFailureSet, FABL.signSamplesF₂Equiv, FABL.uniformPMF}
\difficulty{5}
Let $n$ and $\ell$ be nonnegative integers, and let $\delta$ be rational with
$0<\delta\le 1/2$. Write $N$ for the number of subsets of $[n]$ of size at most $\ell$,
and set $m=2^\ell\bigl(N+\lceil\log_2\lceil 2/\delta\rceil\rceil\bigr)$, where the
logarithmic ceiling is the least integer $k$ such that $2^k$ is at least its argument.
Choose $m$ sign-cube inputs independently and uniformly, and identify each with a vector
in $\mathbb F_2^n$ by the coordinatewise binary-to-sign correspondence. The probability
that evaluation on these inputs fails to be injective on the space of $\mathbb
F_2$-coefficient polynomials spanned by the squarefree monomials of degree at most
$\ell$ is at most $\delta$.
\end{theorem}

\begin{theorem}[Uniqueness from separating samples]
\lean{FABL.coefficients_eq_of_separates_of_satisfies}
\label{FABL.coefficients_eq_of_separates_of_satisfies}
\leanok
\uses{FABL.F₂BooleanFunction, FABL.F₂Cube, FABL.F₂SatisfiesRows, FABL.LowDegreeF₂Coefficients, FABL.SeparatesLowDegreeF₂Coefficients, FABL.functionAlgebraicDegree, FABL.lowDegreeF₂CoefficientsOfFunction, FABL.lowDegreeF₂Eval, FABL.lowDegreeF₂Eval_coefficientsOfFunction, FABL.lowDegreeF₂SampleEquation, FABL.lowDegreeF₂SampleEquation_isSatisfied_iff, FABL.lowDegreeF₂SampleRows}
\difficulty{4}
Let $f:\mathbb F_2^n\to\mathbb F_2$ have algebraic degree at most $\ell$, and let
$x_1,\ldots,x_m\in\mathbb F_2^n$ be sample inputs such that evaluation on these inputs
is injective on the coefficient assignments indexed by subsets of size at most $\ell$.
Any such coefficient assignment that gives the value $f(x_i)$ at every sample input
equals the assignment of the algebraic normal form coefficients of $f$ on those subsets.
\end{theorem}

\begin{definition}[Low-degree coefficient hypothesis]
\lean{FABL.LowDegreeF₂Hypothesis}
\leanok
\uses{FABL.LowDegreeF₂Coefficients}
\statementsource{Analysis-of-Boolean-Functions-by-Ryan-ODonnell}{
  pdf-ffbe9c7b3306-p149-b001,
  pdf-ffbe9c7b3306-p149-b004,
  pdf-ffbe9c7b3306-p167-b005
}
For nonnegative integers $n$ and $\ell$, a low-degree coefficient hypothesis is a
package consisting of a coefficient array indexed by subsets of $[n]$ of size at most
$\ell$, with values in $\mathbb F_2$.
\end{definition}

\begin{definition}[Evaluation of a coefficient hypothesis]
\lean{FABL.LowDegreeF₂Hypothesis.evaluate}
\leanok
\uses{FABL.F₂BooleanFunction, FABL.LowDegreeF₂Hypothesis, FABL.lowDegreeF₂Eval}
\statementsource{Analysis-of-Boolean-Functions-by-Ryan-ODonnell}{
  pdf-ffbe9c7b3306-p149-b001,
  pdf-ffbe9c7b3306-p149-b004,
  pdf-ffbe9c7b3306-p167-b005
}
Let $n$ and $\ell$ be nonnegative integers, and let a low-degree coefficient hypothesis
be given. Its evaluation is the Boolean function on $\mathbb F_2^n$ obtained by
evaluating its coefficient array.
\end{definition}

\begin{definition}[Low-degree sample solver]
\lean{FABL.solveLowDegreeF₂Samples}
\leanok
\uses{FABL.F₂Cube, FABL.LowDegreeF₂Hypothesis, FABL.lowDegreeF₂SampleRows, FABL.solveF₂Rows, FABL.𝔽₂}
\statementsource{Analysis-of-Boolean-Functions-by-Ryan-ODonnell}{pdf-ffbe9c7b3306-p167-b005}
For natural numbers $n,\ell,m$ and $m$ labeled samples consisting of an input in
$\mathbb{F}_2^n$ and a label in $\mathbb{F}_2$, define the returned low-degree
hypothesis by applying the binary linear row solver to the equations formed from those
samples, with one coefficient for each monomial of degree at most $\ell$.
\end{definition}

\begin{theorem}[Consistency of the sample solver]
\lean{FABL.solveLowDegreeF₂Samples_satisfies_of_exists}
\leanok
\uses{FABL.F₂Cube, FABL.F₂SatisfiesRows, FABL.LowDegreeF₂Coefficients, FABL.lowDegreeF₂SampleRows, FABL.solveF₂Rows_satisfies_of_exists, FABL.solveLowDegreeF₂Samples, FABL.𝔽₂}
\difficulty{1}
Let $n,\ell,m$ be natural numbers and choose $m$ samples in $\{0,1\}^n\times\mathbb
F_2$. If some degree-at-most-$\ell$ coefficient family satisfies every equation
determined by these samples, then the coefficient family returned by the specified
solver satisfies every one of those equations.
\end{theorem}

\begin{theorem}[Recovery of a low-degree Boolean function]
\lean{FABL.solveLowDegreeF₂Samples_evaluate_eq}
\leanok
\uses{FABL.F₂BooleanFunction, FABL.F₂Cube, FABL.LowDegreeF₂Hypothesis.evaluate, FABL.SeparatesLowDegreeF₂Coefficients, FABL.coefficientsOfFunction_satisfies_lowDegreeF₂SampleRows, FABL.coefficients_eq_of_separates_of_satisfies, FABL.functionAlgebraicDegree, FABL.lowDegreeF₂CoefficientsOfFunction, FABL.lowDegreeF₂Eval, FABL.lowDegreeF₂Eval_coefficientsOfFunction, FABL.solveLowDegreeF₂Samples, FABL.solveLowDegreeF₂Samples_satisfies_of_exists}
\difficulty{4}
Let $n,\ell,m$ be natural numbers and let $f:\{0,1\}^n\to\mathbb F_2$ have algebraic
degree at most $\ell$. If the $m$ sample inputs separate degree-at-most-$\ell$
coefficient families, then the function represented by the solver's output on the
samples $(x,f(x))$ equals $f$ on all inputs.
\end{theorem}

\begin{definition}[Binary encoding of a labeled sample]
\lean{FABL.binaryLabeledSample}
\statementsource{Analysis-of-Boolean-Functions-by-Ryan-ODonnell}{pdf-ffbe9c7b3306-p167-b005}
\label{FABL.binaryLabeledSample}
\leanok
\uses{FABL.F₂Cube, FABL.Sign, FABL.binaryCubeSignEquiv, FABL.binarySignEquiv, FABL.𝔽₂}
For a natural number $n$, convert a labeled sample consisting of a sign vector in
$\{\pm1\}^n$ and a sign label into a pair consisting of its corresponding vector in
$\mathbb{F}_2^n$ and its corresponding bit in $\mathbb{F}_2$.
\end{definition}

\begin{definition}[Labeled low-degree learner output]
\lean{FABL.lowDegreeF₂PolynomialLearnerLabeledOutput}
\leanok
\uses{FABL.LowDegreeF₂Hypothesis, FABL.Sign, FABL.binaryLabeledSample, FABL.solveLowDegreeF₂Samples}
\statementsource{Analysis-of-Boolean-Functions-by-Ryan-ODonnell}{pdf-ffbe9c7b3306-p167-b005}
For natural numbers $n,\ell,m$ and $m$ labeled samples from $\{\pm1\}^n$ with sign
labels, convert each sample to its binary encoding and apply the low-degree sample
solver. The resulting output is a coefficient assignment over $\mathbb{F}_2$ indexed by
subsets of $\{1,\ldots,n\}$ of size at most $\ell$.
\end{definition}

\begin{theorem}[Binary encoding of a labeled sample]
\lean{FABL.binaryLabeledSample_target}
\label{FABL.binaryLabeledSample_target}
\leanok
\uses{FABL.BooleanFunction, FABL.binaryCubeSignEquiv, FABL.binaryLabeledSample, FABL.booleanFunctionF₂Encoding}
\difficulty{4}
Let $n$ be a natural number, let $f$ map the sign cube of dimension $n$ to the two signs
$-1$ and $1$, and let $x$ be a point of that sign cube. Encoding the sample $(x,f(x))$
as a binary input and binary label gives the pair consisting of the binary encoding of
$x$ and the binary encoding of $f(x)$ evaluated at that encoded input.
\end{theorem}

\begin{theorem}[Exact recovery from separating samples]
\lean{FABL.lowDegreeF₂PolynomialLearnerLabeledOutput_evaluate_eq}
\leanok
\uses{FABL.BooleanFunction, FABL.LowDegreeF₂Hypothesis.evaluate, FABL.SeparatesLowDegreeF₂Coefficients, FABL.binaryCubeSignEquiv, FABL.binaryLabeledSample_target, FABL.booleanFunctionF₂Encoding, FABL.functionAlgebraicDegree, FABL.lowDegreeF₂PolynomialLearnerLabeledOutput, FABL.solveLowDegreeF₂Samples_evaluate_eq}
\difficulty{3}
Let $n,\ell,m$ be natural numbers, and let $f:\{\pm1\}^n\to\{\pm1\}$ have algebraic
degree at most $\ell$ after encoding its inputs and outputs in $\mathbb F_2$. If the $m$
encoded sample inputs distinguish all coefficient families on monomials of degree at
most $\ell$, then the low-degree polynomial learner trained on those inputs with labels
from $f$ evaluates to the encoded function $f$ on every input.
\end{theorem}

\begin{definition}[Work count for a low-degree polynomial learner]
\lean{FABL.lowDegreeF₂PolynomialLearnerWork}
\leanok
\uses{FABL.F₂LinearEquation, FABL.LowDegreeMonomial, FABL.Sign, FABL.binaryLabeledSample, FABL.eliminateF₂Rows, FABL.f₂LinearCoordinateList, FABL.lowDegreeF₂MonomialCount, FABL.lowDegreeF₂SampleRows}
\statementsource{Analysis-of-Boolean-Functions-by-Ryan-ODonnell}{pdf-ffbe9c7b3306-p167-b005}
For natural numbers $n$, $\ell$, and $m$, let each of $m$ samples consist of a point of
the Boolean cube $\{\pm1\}^n$ and a label in $\{\pm1\}$. Form the low-degree linear
equations over $\mathbb F_2$ from these labeled samples, and let $W$ and $P$ be,
respectively, the work count and pivot list produced by eliminating that system; if $M$
is the number of monomials of degree at most $\ell$, the learner's work count is
$m(M+1)+W+|P|(M+1)$.
\end{definition}

\begin{theorem}[Upper bound on learner work]
\lean{FABL.lowDegreeF₂PolynomialLearnerWork_le}
\leanok
\uses{FABL.F₂LinearEquation, FABL.LowDegreeMonomial, FABL.Sign, FABL.binaryLabeledSample, FABL.eliminateF₂Rows, FABL.eliminateF₂Rows_pivots_length_le, FABL.eliminateF₂Rows_work_le, FABL.f₂LinearCoordinateList, FABL.lowDegreeF₂MonomialCount, FABL.lowDegreeF₂PolynomialLearnerWork, FABL.lowDegreeF₂SampleRows}
\difficulty{4}
For any natural numbers $n,\ell,m$ and any list of $m$ labeled examples from
$\{\pm1\}^n$, the learner's work is at most $m(N+1)+Nm(N+2)+N(N+1)$, where $N$ is the
number of subsets of $\{1,\ldots,n\}$ of size at most $\ell$.
\end{theorem}

\begin{definition}[Cubic-scale learner parameter]
\lean{FABL.lowDegreeF₂PolynomialLearnerCubicScale}
\leanok
\uses{FABL.PositiveLearningParameter, FABL.fourierEstimatorFailureBits}
\statementsource{Analysis-of-Boolean-Functions-by-Ryan-ODonnell}{pdf-ffbe9c7b3306-p167-b005}
For natural numbers $n,\ell$ and a rational parameter $\delta$ with $0<\delta\leq 1/2$,
define the scale to be $2^\ell\big((\ell+1)(n+1)^\ell + b(\delta)+2\big)$, where
$b(\delta)$ is the failure-bit count associated with the Fourier estimator.
\end{definition}

\begin{theorem}[Cubic work bound for the learner]
\lean{FABL.scheduledLowDegreeF₂PolynomialLearnerWork_le_cubicScale}
\leanok
\uses{FABL.PositiveLearningParameter, FABL.Sign, FABL.fourierEstimatorFailureBits, FABL.lowDegreeF₂LearningSampleCount, FABL.lowDegreeF₂MonomialCount, FABL.lowDegreeF₂MonomialCount_le, FABL.lowDegreeF₂PolynomialLearnerCubicScale, FABL.lowDegreeF₂PolynomialLearnerWork, FABL.lowDegreeF₂PolynomialLearnerWork_le}
\difficulty{5}
For nonnegative integers $n,\ell$, a rational parameter $\delta$ with $0<\delta\le 1/2$,
and any labeled sample sequence of the scheduled sample count, the low degree polynomial
learner's work is at most $3K^3$, where
$K=2^\ell\bigl((\ell+1)(n+1)^\ell+b(\delta)+2\bigr)$ and
$b(\delta)=\lceil\log_2\lceil2/\delta\rceil\rceil$.
\end{theorem}

\begin{theorem}[Fixed parameter work envelope]
\lean{FABL.scheduledLowDegreeF₂PolynomialLearnerWork_le_fixedParameterEnvelope}
\leanok
\uses{FABL.PositiveLearningParameter, FABL.Sign, FABL.fourierEstimatorFailureBits, FABL.lowDegreeF₂LearningSampleCount, FABL.lowDegreeF₂PolynomialLearnerCubicScale, FABL.lowDegreeF₂PolynomialLearnerWork, FABL.scheduledLowDegreeF₂PolynomialLearnerWork_le_cubicScale}
\difficulty{4}
For nonnegative integers $n,\ell$, a rational parameter $\delta$ with $0<\delta\le 1/2$,
and any labeled sample sequence of the scheduled sample count, the low degree polynomial
learner's work is at most $3\bigl(2^\ell(\ell+b(\delta)+3)\bigr)^3(n+1)^{3\ell}$, where
$b(\delta)=\lceil\log_2\lceil2/\delta\rceil\rceil$.
\end{theorem}

\begin{definition}[Low-degree polynomial learner program]
\lean{FABL.lowDegreeF₂PolynomialLearnerProgram}
\leanok
\uses{FABL.LearningProgram, FABL.LowDegreeF₂Hypothesis, FABL.lowDegreeF₂PolynomialLearnerLabeledOutput, FABL.lowDegreeF₂PolynomialLearnerWork}
\statementsource{Analysis-of-Boolean-Functions-by-Ryan-ODonnell}{pdf-ffbe9c7b3306-p167-b005}
For natural numbers $n$, $\ell$, and $m$, define a random-example learning program that
draws exactly $m$ independent uniform labeled examples from the sign cube $\{\pm1\}^n$,
charges the local work specified by the low-degree learner for those samples, and
returns the learner's resulting hypothesis, given by a coefficient in $\mathbb{F}_2$ for
each monomial of degree at most $\ell$.
\end{definition}

\begin{theorem}[Low-degree learner run distribution]
\lean{FABL.runWithCost_lowDegreeF₂PolynomialLearnerProgram}
\leanok
\uses{FABL.BooleanFunction, FABL.LearningProgram.runWithCost, FABL.lowDegreeF₂PolynomialLearnerLabeledOutput, FABL.lowDegreeF₂PolynomialLearnerProgram, FABL.lowDegreeF₂PolynomialLearnerWork, FABL.uniformPMF}
\difficulty{3}
For natural numbers $n$, $\ell$, and $m$, and a Boolean function
$f:\{−1,1\}^n\to\{−1,1\}$, running the low-degree polynomial learner with target $f$ has
the same probability law as drawing a uniformly random sequence $(x_i)_{i=1}^m$ from
$\{−1,1\}^n$, returning the learner’s output on the labeled samples $(x_i,f(x_i))$, and
assigning cost $(m,0,m+W)$, where $W$ is the learner’s charged work on those samples.
\end{theorem}

\begin{theorem}[Oracle cost of the learner]
\lean{FABL.lowDegreeF₂PolynomialLearnerProgram_oracleCost}
\leanok
\uses{FABL.BooleanFunction, FABL.LearningCost, FABL.LearningProgram.runWithCost, FABL.LowDegreeF₂Hypothesis, FABL.lowDegreeF₂PolynomialLearnerProgram, FABL.runWithCost_lowDegreeF₂PolynomialLearnerProgram}
\difficulty{4}
For any natural numbers $n,\ell,m$, target function on $\{\pm1\}^n$, and outcome in the
support of running the low-degree polynomial learner with $m$ random examples on that
target, the outcome's random-example cost is $m$ and its membership-query cost is zero.
\end{theorem}

\begin{definition}[Scheduled low-degree polynomial learner]
\lean{FABL.scheduledLowDegreeF₂PolynomialLearnerProgram}
\leanok
\uses{FABL.LearningProgram, FABL.LowDegreeF₂Hypothesis, FABL.PositiveLearningParameter, FABL.lowDegreeF₂LearningSampleCount, FABL.lowDegreeF₂PolynomialLearnerProgram}
\statementsource{Analysis-of-Boolean-Functions-by-Ryan-ODonnell}{pdf-ffbe9c7b3306-p167-b005}
For natural numbers $n$ and $\ell$, and a rational parameter $\delta$ with $0<\delta\leq
1/2$, define a random-example learning program by running the low-degree $\mathbb{F}_2$
polynomial learner with the sample count prescribed by $n$, $\ell$, and $\delta$. Its
output is a hypothesis assigning an $\mathbb{F}_2$ coefficient to each subset of
$\{1,\ldots,n\}$ of size at most $\ell$.
\end{definition}

\begin{theorem}[Learner failure implies separation failure]
\lean{FABL.lowDegreeF₂PolynomialLearner_failure_subset_separationFailure}
\leanok
\uses{FABL.BooleanFunction, FABL.LowDegreeF₂Hypothesis.evaluate, FABL.SeparatesLowDegreeF₂Coefficients, FABL.booleanFunctionF₂Encoding, FABL.functionAlgebraicDegree, FABL.lowDegreeF₂PolynomialLearnerLabeledOutput, FABL.lowDegreeF₂PolynomialLearnerLabeledOutput_evaluate_eq, FABL.lowDegreeF₂SeparationFailureSet, FABL.signSampleLowDegreeF₂SeparationFailureSet, FABL.signSamplesF₂Equiv}
\difficulty{4}
Let $n,\ell,m$ be natural numbers, and let a target function on $\{\pm1\}^n$ have
encoded algebraic degree at most $\ell$. The set of sample-input lists on which the
learner, trained with labels from the target, fails to evaluate to the encoded target is
contained in the set of lists whose encoded inputs do not distinguish all low-degree
coefficient families.
\end{theorem}

\begin{theorem}[Scheduled low degree f₂polynomial learner program failure probability le]
\lean{FABL.scheduledLowDegreeF₂PolynomialLearnerProgram_failureProbability_le}
\leanok
\uses{FABL.BooleanFunction, FABL.LearningProgram.eventProbability, FABL.LowDegreeF₂Hypothesis.evaluate, FABL.PositiveLearningParameter, FABL.booleanFunctionF₂Encoding, FABL.functionAlgebraicDegree, FABL.lowDegreeF₂LearningSampleCount, FABL.lowDegreeF₂PolynomialLearnerLabeledOutput, FABL.lowDegreeF₂PolynomialLearner_failure_subset_separationFailure, FABL.measure_signSampleLowDegreeF₂SeparationFailureSet_scheduled_le, FABL.runWithCost_lowDegreeF₂PolynomialLearnerProgram, FABL.scheduledLowDegreeF₂PolynomialLearnerProgram, FABL.uniformPMF}
\statementsource{Analysis-of-Boolean-Functions-by-Ryan-ODonnell}{pdf-ffbe9c7b3306-p167-b005}
\difficulty{5}
% dataset-pipeline: stub -- prose pending build_informal_statements.py
Statement pending regeneration.
\end{theorem}

\begin{theorem}[Low degree learner sample count]
\lean{FABL.scheduledLowDegreeF₂PolynomialLearnerProgram_sampleCount}
\leanok
\uses{FABL.PositiveLearningParameter, FABL.lowDegreeF₂LearningSampleCount, FABL.lowDegreeF₂MonomialCount}
\difficulty{1}
For nonnegative integers $n$ and $\ell$ and a rational parameter $\delta$ with
$0<\delta\le 1/2$, the scheduled low degree polynomial learner's sample count is
$2^\ell\bigl(M(n,\ell)+\lceil\log_2\lceil 2/\delta\rceil\rceil\bigr)$, where $M(n,\ell)$
is the number of subsets of an $n$-element set having cardinality at most $\ell$.
\end{theorem}
